\introsection{Научная новизна исследования}

Научная новизна проведённого исследования заключается в разработке и практической реализации ряда оригинальных подходов, направленных на повышение качества генерации SQL-запросов в задачах класса Text-to-SQL в условиях сложных и слабо формализованных схем данных.

В рамках работы предложен адаптивный метод семантического поиска контекста схемы базы данных. В отличие от традиционных решений, ограничивающихся использованием только названий таблиц и столбцов, разработанный подход опирается на многоуровневое индексирование метаинформации. В процесс индексирования включаются не только структурные элементы схемы (названия таблиц и колонок), но и дополнительные источники контекста, такие как комментарии к полям, а также репрезентативные примеры данных. Это позволяет формировать более полное и семантически насыщенное представление структуры базы данных. В результате достигается повышение точности генерации SQL-запросов, особенно в условиях работы со сложными схемами, содержащими большое количество взаимосвязанных таблиц (порядка 5 и более).

Кроме того, предложен оригинальный алгоритм многошаговой самокоррекции сгенерированных запросов (Self-Correction), реализующий концепцию взаимодействия «критика» и «исполнителя». В рамках данного подхода интеллектуальный агент после генерации SQL-запроса выполняет его в системе управления базами данных и анализирует возникающие ошибки. Полученные диагностические сообщения СУБД интерпретируются как источник обратной связи, на основе которого агент итеративно корректирует запрос, устраняя как синтаксические, так и логические ошибки. Такой механизм позволяет существенно повысить надёжность системы, обеспечивая автоматическое исправление ошибок без участия пользователя и снижая зависимость от первоначального качества генерации.
